% ==================================================================
% id:       rv-003
% topic:    משתנים מקריים
% subtopic: משתנים מקריים מוכרים
% solution: full
% verified: no
% ==================================================================
שנה אחרי שנה, פרופסור אקס מלמד את הקורס "מבוא למוטנטים". הקורס נחשב קשה, וידוע שההסתברות שלפחות סטודנט אחד יקבל 100 בקורס היא $0.02$, בכל שנה באופן בלתי תלוי בשנים הקודמות. פתרו כל אחד מהסעיפים הבאים בנפרד.
\begin{solbox}
\unverified
נסמן ב-$X$ את מספר השנים עד שסטודנט קיבל 100 בקורס. כלומר $X\sim\Geo(0.02)$.
\end{solbox}
\begin{enumerate}
\item מה ההסתברות שרק בשנה הרביעית היה, לראשונה בתולדות הקורס, סטודנט שקיבל 100?
\begin{solbox}
\unverified
$\PP(X=4)=(1-0.02)^3\cdot0.02=0.98^3\cdot0.02\approx\ans{0.0188}$.
\end{solbox}
\item מה ההסתברות שבשש השנים הראשונות לא היה אף סטודנט שקיבל 100?
\begin{solbox}
\unverified
6 שנים ללא "הצלחה": $\PP(X>6)=0.98^6\approx\ans{0.8858}$.
\end{solbox}
\item מה ההסתברות שבעשר השנים הראשונות היו בדיוק שנתיים שבהן סטודנט קיבל 100?
\begin{solbox}
\unverified
כאן סופרים הצלחות במשך 10 שנים, כלומר זה משתנה בינומי. נסמן ב-$Y$ את מספר השנים, מתוך 10 הראשונות, שבהן סטודנט קיבל 100. אז $Y\sim\Bin(10,0.02)$ ו-
\[
\PP(Y=2)=\binom{10}2(0.02)^2(0.98)^8\approx\ans{0.0153}.
\]
\end{solbox}
\item בכל שנה, פרופ' מקגונגל (עמיתה של פרופ' אקס) מלמדת במקביל את הקורס "מבוא לקסמים". ההסתברות שלפחות סטודנט אחד יקבל 100 בקורס שלה היא $0.05$, באופן בלתי תלוי בשנים הקודמות. בנוסף, הקורסים של שני הפרופסורים בלתי תלויים זה בזה.
שני הפרופסורים הסכימו שבשנה הראשונה שבה סטודנט מקבל 100 בקורס שלהם, הם יוצאים לשנת שבתון.
\begin{enumerate}
\item מה ההסתברות ששניהם יצאו לשבתון אחרי שנה אחת בלבד?
\item מה ההסתברות ששניהם יצאו לשבתון בדיוק באותה שנה?
\end{enumerate}
\begin{solbox}
\unverified
נגדיר $Z$ = מספר השנים עד שסטודנט קיבל 100 בקורס של פרופ' מקגונגל. אז $Z\sim\Geo(0.05)$, ו-$X,Z$ בלתי תלויים.
\begin{enumerate}
\item $\PP(X=1,Z=1)=\PP(X=1)\PP(Z=1)=0.02\cdot0.05=\ans{0.001}$.
\item נסמן $A_n=\set{X=n,\ Z=n}$, המאורע ששניהם יצאו לשבתון בשנה ה-$n$. המאורע המבוקש הוא $A=\bigcup_{n\ge1}A_n$, והמאורעות $A_n$ זרים בזוגות. לפי אי-התלות:
\[
\PP(A_n)=\PP(X=n)\PP(Z=n)=0.98^{n-1}\cdot0.02\cdot0.95^{n-1}\cdot0.05=0.001\cdot q^{n-1},
\]
כאשר $q=0.98\cdot0.95$. לכן, לפי סכום של טור הנדסי:
\[
\PP(A)=\sum_{n=1}^\infty\PP(A_n)=0.001\sum_{n=1}^\infty q^{n-1}=\frac{0.001}{1-q}\approx\ans{0.0145}.
\]

\end{enumerate}
\end{solbox}
\end{enumerate}
